% !TEX root = ../main.tex
\chapter{要求が届くまで}
\label{chap:request-delivery}

\begin{chapterabstract}
前章のIPは、パケットを宛先の端末付近まで運びました。
この章では、端末内のアプリケーションを識別し、必要な信頼性と安全性を加えます。
最後にHTTP要求をサーバーへ渡し、言語モデルを呼び出す直前まで進みます。
\end{chapterabstract}

\begin{learninggoals}
  \item TCP、UDP、QUICの機能と保証範囲を説明できます。
  \item TLSが機密性、完全性、認証をどのように支えるかを説明できます。
  \item URL、HTTP、キャッシュ、プロキシ、CDNの役割を要求の順序へ配置できます。
\end{learninggoals}

\chapterimage{ch06-delivery.png}{HTTPSで接続を作り、キャッシュの確認後にサーバーへ進む場合の概念図}{fig:ch06-delivery}

\section{ポート番号}

\term{ポート番号}は、同じIPアドレス上の通信の受け口を区別する番号です。
TCPやUDPなど各トランスポートプロトコルが持つ16ビットの番号で、0から65535までを表せます。
\term{ソケット}は、通信端点を表すOSの抽象化です。
典型的には、プロトコル、送信元と宛先のIPアドレス、送信元と宛先のポート番号で通信を識別します。
ポート番号だけで利用者やアプリケーションの正当性を認証するわけではありません。

\section{TCP接続}

\term{TCP}はTransmission Control Protocolの略で、順序のあるバイト列を届けるトランスポート層のプロトコルです。
\term{コネクション}は、両端が通信状態を保持している論理的な関係です。
\term{3ウェイハンドシェイク}は、TCP接続を始めるための3段階の交換です。

\term{SYN}は、接続開始と初期シーケンス番号を伝える制御情報です。
\term{ACK}は、どこまで受信したかを確認する制御情報です。
\term{FIN}は、送信方向を終了する意思を伝える制御情報です。
接続開始側がSYNを送り、受信側がSYNとACKを返し、開始側がACKを返す順序を、\code{SYN → SYN-ACK → ACK}と表します。これにより、両方向の初期シーケンス番号と受信能力を確認して接続を確立します。
接続の確立は、HTTP処理の成功や相手アプリケーションの正しさまでは保証しません。

\section{TCPのデータ転送}

\term{シーケンス番号}は、バイト列の位置を表す値です。
\term{順序制御}は、受信したデータを元の順序へ並べます。
\term{再送制御}は、確認応答が得られないデータを再び送ります。
\term{タイムアウト}は、応答を待ち続けず再送や失敗判断へ移るための時間条件です。

前章の\term{RTT}は、確認応答を待つ時間や再送判断に影響します。
\term{パケット損失}が起きると、TCPは不足部分を補うまで上位層へ渡せない場合があります。
\source{\href{https://datatracker.ietf.org/doc/html/rfc9293}{RFC 9293: Transmission Control Protocol}}

図\ref{fig:06-tcp-exchange}では、100バイトずつ送る例を使います。
ACK 200は、次に200番のバイトを受け取りたいという確認応答です。
300〜399番が先に届いても、200〜299番が不足している間はACK 200を返します。
不足部分を再送で補うと、400番まで進めます。
実際の再送開始は、再送タイマーや損失検出の規則で決まります。
\sectionimage{generated/06-tcp-exchange.png}{TCP接続の開始と、順序をそろえる再送の例}{fig:06-tcp-exchange}


\section{フロー制御と輻輳制御}

\term{フロー制御}は、受信側の処理能力に合わせて送信量を調整します。
\term{輻輳制御}は、網の混み具合に合わせて送信量を調整します。
\term{輻輳ウィンドウ}は、輻輳制御が未確認のまま送れる量を制限する値です。
フロー制御は受信側を守り、輻輳制御は網の混雑を抑えます。

\section{UDP}

\term{UDP}はUser Datagram Protocolの略で、宛先ポートへデータグラムを運ぶ薄いプロトコルです。
UDPは接続状態、順序制御、再送制御を標準では提供しません。
必要な性質は、アプリケーションまたはUDP上の別プロトコルが補います。
「UDPは常に高速」という意味ではなく、用途に合わせて制御を構成しやすいことが特徴です。

\section{QUIC}

\term{QUIC}は、UDP上で接続管理、暗号化、再送、多重化を実現するトランスポートプロトコルです。
\term{多重化}は、複数の論理的なデータ流を一つの接続上で並行に扱うことです。
\term{HTTP/3}は、HTTPの意味論をQUIC上で運びます。
\term{TLS 1.3}の暗号学的な仕組みはQUICの接続確立へ統合されています。

\term{ヘッドオブラインブロッキング}は、先頭側の不足によって後続の処理まで待たされる現象です。
QUICはストリーム間の待ちを分離しますが、同じストリーム内の順序待ちやネットワーク遅延は残ります。
\source{\href{https://datatracker.ietf.org/doc/html/rfc9000}{RFC 9000: QUIC}}

\section{HTTP}

\term{HTTP}はHypertext Transfer Protocolの略で、要求と応答の意味を定めるアプリケーション層の約束です。
\term{ヘッダー}は、内容の種類、キャッシュ条件、認証情報などのメタデータを運びます。
\term{ボディ}は、要求または応答の本体を運びます。
HTTPはメッセージの意味を定めますが、業務処理が正しいことまでは保証しません。

\section{HTTPメソッドとステータス}

\term{HTTPメソッド}は、対象へ行いたい操作の意味を表します。
\term{GET}は、対象の表現を取得するために使います。
\term{POST}は、サーバー側へ処理を依頼するときに使います。
\term{PUT}は、指定した対象の状態を作成または置換する意味を持ちます。
\term{DELETE}は、指定した対象の削除を依頼する意味を持ちます。

\term{ステータスコード}は、応答結果を3桁の分類で示します。
2xxは成功、4xxはクライアント側の要求に関する問題、5xxはサーバー側の処理に関する問題を表します。
個別の意味はRFC 9110の定義を確認します。
\source{\href{https://datatracker.ietf.org/doc/html/rfc9110}{RFC 9110: HTTP Semantics}}

\section{TLS}

\term{TLS}はTransport Layer Securityの略で、通信を暗号学的に保護するプロトコルです。
\term{TLSハンドシェイク}は、保護された通信を始めるための交換手順です。
暗号パラメータを交渉し、通信を保護する鍵を確立します。
証明書を使う方式では、提示された証明書と署名などを確認して相手を認証します。
\term{公開鍵}は、対応する秘密鍵を公開せずに検証や鍵共有へ使える鍵です。
\term{秘密鍵}は、所有者が秘密に管理する鍵です。
\term{共通鍵}は、暗号化と復号で共有する鍵です。

\term{暗号化}は、鍵を持たない第三者が内容を読みにくくする処理です。
\term{完全性}の保護は、通信途中の改ざんや破損を検出する仕組みです。
\term{認証}は、通信相手が想定した主体であることを確認する仕組みです。
TLSはサーバーのアプリケーション処理や返答内容の正しさを保証しません。
\source{\href{https://datatracker.ietf.org/doc/html/rfc8446}{RFC 8446: TLS 1.3}}

\section{証明書検証}

\term{証明書}は、公開鍵と主体名などを署名によって結び付けたデータです。
\term{認証局}は、証明書を発行または署名する信頼主体です。
認証局はCertificate Authorityを略して\term{CA}とも呼びます。
\term{証明書チェーン}は、サーバー証明書から信頼済みのルートまで署名関係をたどる並びです。

クライアントは、署名、期限、ホスト名、信頼経路、失効情報を利用可能な範囲で検証します。
警告を無視して接続すると、相手を確認する前提が崩れます。

\section{HTTPの各バージョン}

\term{HTTP/1.1}は、テキスト形式のメッセージをTCP上で運ぶ方式です。
\term{HTTP/2}は、バイナリフレームとストリーム多重化をTCP上で使います。
HTTP/2では複数ストリームを扱えますが、TCP層の欠損が接続全体へ影響する場合があります。
\term{HTTP/3}は、QUICのストリーム上でHTTPの意味論を運びます。
バージョンが違っても、メソッドやステータスの基本的な意味は共通です。

\section{URLから接続まで}

\term{URI}はUniform Resource Identifierの略で、資源を識別する文字列です。
\term{URL}はUniform Resource Locatorの略で、資源の場所を示すURIの一種です。
\term{スキーム}は、\code{https}のように利用する方式を示します。
\term{ホスト}は、接続先を示すドメイン名またはIPアドレスです。
\term{パス}は、ホスト内の対象位置を示します。
\term{クエリ}は、\code{?}以降へ追加条件を表します。

ブラウザはURLを分解し、DNSでホストを解決します。
次にTCPまたはQUICで通信を始め、TLSで相手と通信内容を保護します。
最後にHTTP要求を送ります。

\section{サーバーの手前}

\term{キャッシュ}は、以前の応答を条件付きで再利用して処理や転送を減らします。
\term{プロキシ}は、クライアントとサーバーの間で要求や応答を中継します。
\term{CDN}はContent Delivery Networkの略で、利用者に近い拠点からコンテンツを配信します。
\term{Service Worker}は、ブラウザ内でネットワーク要求を仲介できるスクリプトです。

\term{fetch}は、ブラウザのJavaScriptからHTTP要求を開始するAPIです。
これらの中継や再利用があるため、観測した応答が必ず元サーバーで生成されたとは限りません。

\section{手を動かす：通信の時間}

\term{curl}は、URLで指定した相手へ通信するコマンドラインツールです。
\term{DevTools Network}は、ブラウザで要求の順序、ヘッダー、状態、時間を観測する画面です。
許可された対象に対して次の例を実行します。

\begin{lstlisting}[language=bash]
curl -o /dev/null -sS \
  -w 'dns=%{time_namelookup} connect=%{time_connect} tls=%{time_appconnect} ttfb=%{time_starttransfer} total=%{time_total}\n' \
  https://example.com/
\end{lstlisting}

各値は、開始からその段階までの累積秒数です。
新しいTCP接続でHTTPSを使い、リダイレクトやプロキシを使わない場合を考えます。
接続部分の目安は\code{connect - dns}です。
TLS部分の目安は\code{tls - connect}です。
\code{ttfb - tls}には、要求の転送、サーバー処理、応答が戻る時間が含まれます。
同じ条件で複数回測り、時刻と回数をそろえて比較します。
\source{\href{https://curl.se/docs/manpage.html}{curl manual, write-out: time\_namelookup / time\_connect / time\_appconnect / time\_starttransfer}}

\section{サーバーの内側へ}

HTTP要求は、OSのソケットからサーバープロセスへ渡ります。
アプリケーションは注文番号を検証し、注文データを取得し、言語モデルへ渡す事実を組み立てます。
次章では、その入力が数値表現へ変わり、応答文として生成されるまでを扱います。

\section{旅をたどり直す}

注文確認AIのURLは、接続方式、ホスト、パスを示します。
DNSとIPが宛先への基盤を用意し、TCPまたはQUICが通信状態を管理します。
TLSが通信相手と内容を保護し、HTTPが要求の意味をサーバーへ伝えます。
\takeaway{要求は、端末内の識別、転送、接続、暗号学的保護を経て届きます。HTTPが、その要求の意味を相手へ伝えます。}

\subsection*{章末の案内語}

この表は、第6章で説明した語のうち、端末内の届け分けから保護されたHTTP要求までをたどる語に絞っています。

\noindent\begin{tabularx}{\linewidth}{@{}>{\bfseries\raggedright\arraybackslash\hyphenpenalty=10000}p{32mm}>{\raggedright\arraybackslash}p{35mm}Y@{}}
  \toprule
  この章の語 & 最初に説明した節 & 一言で戻る \\
  \midrule
  \guideterm{ポート番号} & ポート番号 & プロトコルごとの16ビット名前空間で通信端点を識別します。 \\
  \guideterm{TCP} & TCP接続 & 順序のあるバイト列を接続状態とともに届けます。 \\
  \guideterm{3ウェイハンドシェイク} & TCP接続 & SYN、SYN-ACK、ACKで接続を確立します。 \\
  \guideterm{UDP} & UDP & 接続や再送を標準では持たないデータグラム配送です。 \\
  \guideterm{QUIC} & QUIC & UDP上で暗号化、再送、多重化を提供します。 \\
  \guideterm{HTTP} & HTTP & 要求と応答の意味を定めます。 \\
  \guideterm{TLSハンドシェイク} & TLS & 暗号パラメータを交渉し、鍵を確立し、相手を認証します。 \\
  \guideterm{証明書} & 証明書検証 & 公開鍵と主体名などを署名で結び付けます。 \\
  \guideterm{URL} & URLから接続まで & 接続方式、ホスト、対象位置などを一つの文字列で示します。 \\
  \guideterm{キャッシュ} & サーバーの手前 & 条件を満たす以前の応答を再利用します。 \\
  \bottomrule
\end{tabularx}

